\sscp{Issues with languages that have more than five vowels}{02-04-06}
Across the \tsc{\ili{Bima}-Lembata} subgroup of Austronesian languages,
and also in the Papuan \tsc{\ili{Alor}-Pantar} languages,
in addition to phonemic contrasts between phonetically short
(single) and long (double) vowels,
there are other widespread patterns that involve schwa /ə/,
reduced vowels (of shortened duration),
centralised vowels (where the five vowels move slightly towards
schwa /ə/), de-stressing of penultimate syllables with schwa /ə/,
or sometimes a stressed penultimate schwa /ə/,
and phonetic lengthening of the following intervocalic consonant.
In several languages there is an interaction between several of these phenomena.

There are a variety of strategies found in various sources
for the transcription of these different and very subtle phenomena.
This is further complicated as Indonesian has a phonemic
six-vowel system with schwa /ə/,
but /ə/ and /e/ are orthographically
under-differentiated and both are represented by ‹e›
\citep{Grimes1996}.\footnote{
		The mid-front vowel /e/ has a very low functional load in Indonesian
		(low frequency), whereas /ə/ has a very high functional load.}
Many native speakers
of local languages, or Indonesian scholars
trying to document languages
in this region often transcribe their data with only the five
vowels of the typewriter or basic keyboard,
and no other indication.

In languages with schwa /ə/,
reduced vowels, or centralised vowels, sometimes nothing is indicated,
even though there is phonemic contrast in identical environments.
Sometimes the investigator marks something on the vowels.
Such marking may or may not be consistent.
Sometimes the investigator writes double consonants to indicate
the phonetic lengthening of the following consonant which is
often triggered by a preceding schwa /ə/,
and is entirely predictable.
However the schwa that triggers the phonetic lengthening is not
indicated by many attempting to write or document those languages.
Such a strategy is somewhat effective for words shaped CəCV → ‹CaCCV›
or ‹CeCCV›, but it does not work for words shaped CəV or əV which have no
following consonant to double.

The \tsc{\ili{Bima}-Lembata} subgroup is sandwiched between other Austronesian
languages to the west
and north, and Papuan \tsc{\ili{Alor}-Pantar} languages to the \ili{east}.
\tsc{PAn} is reconstructed as having a four-vowel system with
*i, *ə, *a, *u \citep[173]{Blust2013}.
\citet{Smith2021} has pointed out that in many Austronesian languages
in the Philippines, western Indonesia, and Sulawesi penultimate schwa /ə/ is oftenː
a) unstressed, b) of short duration,
and/or c) accompanied by phonetic lengthening
of the following consonant.

Besides the patterned differences between the phonetic realisation
of schwa in \tsc{\ili{Bima}-Lembata} languages,
it may be significant that the phonetic realisation of schwa
in these languages is very similar to vocalic phenomena
in nearby Papuan \tsc{\ili{Alor}-Pantar} languages such as \ili{Klon}, \ili{Wersing}, and \ili{Klamu}.

The issue with the transcription of schwa
and reduced vowels in \tsc{\ili{Alor}-Pantar} languages is
seen with alternate language names such as \ili{Klon}, Kelong, Kalong.
\citet[35]{Baird2008} notes that in \ili{Klon} “schwa only occurs in unstressed
syllables.” At a phonology-orthography workshop in November 2019
with a number of native speakers of the \ili{Klon} \ili{Bring} dialect,
it was found that \ili{Klon} has seven phonemic vowels /i ɪ e a o ʊ
u/ with only one central vowel /a/
and contrastive vowel length in all seven vowels.
There is no phonemic schwa,
but there is a non-phonemic epenthetic vowel used to break up consonant clusters.
It is unstressed, of short duration and not syllabic.
It is mostly central,
and often takes on the colour of nearby vowels \citep{Jones2020a}.
Native speakers sometimes refer to
this as a “fast vowel” (Indonesian: \it{vokal cəpat}).
So /mgol/ ‘banana’ → [məgol] {\tl} [mŏgol] {\tl} [mgol],
/ngan/ ‘thing, matter’ → [nəgan] {\tl} [năgan] {\tl} [ngan],
/ngeen/ ‘give us’ → [nəgeːn] {\tl} [nĕgeːn] {\tl} [ngeːn], etc.

Although the Papuan language of \ili{Wersing} spoken on \ili{Alor} has only
a five-vowel base system with long
and short contrasts, it also has a similar phenomenon of epenthetic
schwa-like “fast vowels” used to break up consonant clusters,
as in /ptori/ → [pətori] ‘boiled corn’,
/t-teŋ/ → [təteŋ] ‘our hands’,
/tma/ → [təma] ‘sea’ \citep{Jones2020b}.

The interplay between vowel quality,
vowel length, and consonant length appears to be an areal phenomenon
found in the \tsc{\ili{Bima}-Lembata} nodes of \tsc{Flores-Lembata},
\tsc{Sumba-Havu}, and also in some Papuan \tsc{\ili{Alor}-Pantar} languages.
In addition to lengthening of the following consonant in \tsc{Sumba-Havu}
following a penultimate schwa described in \srf{03-02-01} and \srf{03-03-01},
the languages of western Sumba show a complex interplay between
the quality and length of the penultimate vowel,
and some phonetic lengthening of the following intervocalic consonant.
In some cases all that remains as traces of a certain historical
vowel is the lengthened consonant (see \srf{03-03-03-01}).

\largerpage
\ili{Klamu}, spoken on \ili{Pantar}, manifests most of the phenomena discussed
above in a single language.\footnote{
		There is some preliminary indication that other languages
		of \ili{Pantar}, such as \ili{Teiwa}, may have similar
		complexities to \ili{Klamu} (Amos Sir, pers.
		comm.).
		The “fast vowel” is found across most \ili{Alor}-\ili{Pantar} languages.}
\ili{Klamu} has regular vowels,
long vowels, and short penultimate vowels that trigger phonetic
lengthening of the following consonant.
The five base vowels thus show a three-way contrast in length,
and there is interplay between the short vowels with compensatory
lengthening of the following consonant.
There is also an epenthetic ``fast vowel'' with initial CC clusters
similar to that described for \ili{Wersing} and \ili{Klon}.
The short vowels are distinguished by three phenomena: a) they
are of slightly shorter duration [ă ĕ ĭ ŏ ŭ] than the equivalent
base vowel (V); b) there is vowel reduction in which the short
V are slightly centralised in vowel quality (moving towards schwa
[ə], but only the low vowel /a/ actually reaches a full,
albeit slightly shortened [ə]; c) this shorter duration
and slight centralisation combines with a slight lengthening
[ˑ] of the following C.\footnote{
		The lengthening of the following
		C in \ili{Klamu} is indicated by a half-length [ˑ],
		and to Grimes' ears is not as long as the lengthening of the
		C [ː] following /ə/ in \ili{Havu}
		and \ili{Dhao} (see discussion
		and examples in \srf{03-03-01-02} and \srf{03-03-02-HD}).}
This combined phenomenon
is indicated in the locally derived orthography by a doubling
of medial consonants (including digraphs),
with no marking on the preceding V.
Doubling of the consonants is also found in the writing of family
names, indicating the orthographic strategy has been around for some time.
Word stress is normally penultimate,
including on the reduced V.
Almost all of the hundreds of examples found involve combinations
of /i, a, u/.
Very few involve mid vowels /e, o/.
These phenomena are illustrated in example \qf{02-44}.
(Data from Grimes' fieldnotes.)

{\wg\begin{xltabular}{\linewidth}{>{\hspace{-\tabcolsep}}l<{\hspace{-\tabcolsep}}>{\hspace{-\tabcolsep}\enspace}llX}
\lsptoprule
\tbx{02-44}		&	Orth.	&	Phonetic	&	Gloss		\\ \midrule
\endhead
&	‹wala›	&	[ˈwɑlɑ]	&	‘rock, stone’		\\
	&	‹walla›	&	[ˈwălˑɑ] {\tl} [ˈwəlˑɑ]	&	‘full, packed, crowded’		\\
	&	‹aci›	&	[ˈɑʧi]	&	‘peel skin (of animal or plant)’		\\
	&	‹acci›	&	[ˈăʧˑi] {\tl} [ˈəʧˑi]	&	‘undo knot’		\\
	&	‹baqa›	&	[ˈbɑqɑ]	&	‘hole, opening, orifice’		\\
	&	‹baqqa›	&	[ˈbăqˑɑ] {\tl} [ˈbəqˑɑ]	&	‘body’		\\
	&	‹raaqu›	&	[ˈrɑːqu]	&	‘pick, harvest (grain)’		\\
	&	‹raqqu›	&	[ˈrăqˑu] {\tl} [ˈrəqˑu]	&	‘two’		\\
	&	‹rana›	&	[ˈrɑnɑ]	&	‘face to face’		\\
	&	‹ranna›	&	[ˈrănˑɑ] {\tl} [ˈrənˑɑ]	&	‘how much, how many’		\\
	&	‹lala›	&	[ˈlɑlɑ]	&	‘promise’		\\
	&	‹lalla›	&	[ˈlălˑɑ] {\tl} [ˈləlˑɑ]	&	‘loose’		\\
	&	‹lilla›	&	[ˈlɪlˑɑ] {\tl} [ˈlĭlˑɑ] {\tl} [ˈləlˑɑ]	&	‘sick, hurt’		\\
	&	‹iila›	&	[ˈiːlɑ]	&	‘water’		\\
	&	‹illa›	&	[ˈɪlˑɑ] {\tl} [ˈĭlˑɑ] {\tl} [ˈəlˑɑ]	&	‘fly (v)’		\\
	&	‹banga›	&	[ˈbɑŋɑ]	&	‘ask, request, bargain’		\\
	&	‹bangnga›	&	[ˈbăŋˑɑ] {\tl} [ˈbəŋˑɑ]	&	‘life’		\\
	&	‹wana›	&	[ˈwɑnɑ]	&	‘connect’		\\
	&	‹wanna›	&	[ˈwănˑɑ] {\tl} [ˈwənˑɑ]	&	‘exists’		\\
	&	‹billi›	&	[ˈbɪlˑi] {\tl} [ˈbəlˑi]	&	‘see, check out, look in on’		\\
	&	‹gagga›	&	[ˈgăgˑɑ] {\tl} [ˈgəgˑɑ]	&	‘he, she, it (one of several \tsc{3sg} pronouns)’		\\
	&	‹ecicci›	&	[ɛˈʧɪʧˑi] {\tl} [ɛˈʧəʧˑi]	&	‘true’		\\
	&	‹akku›	&	[ˈăkˑu] {\tl} [ˈəkˑu]	&	‘say, that (complementiser for speech acts)’		\\
	&	‹cukku›	&	[ˈʧŭkˑu] {\tl} [ˈʧəkˑu]	&	‘enter, go into’		\\
	&	‹gsarru›	&	[ˈgsărˑu] {\tl} [ˈgsərˑu]	&	‘see-him (undergoer)’ /g-sarru/		\\ \hhline{~}
	&						&	{\tl} [gəˈsərˑu]	&	\hp{ }		\\
	&	‹tlaqqu›	&	[ˈtlăqˑu] {\tl} [ˈtləqˑu]	&	‘fever and chills,		\\ \hhline{~}
	&						&	{\tl} [təˈləqˑu]		&		\hp{`}symptoms of malaria’	\\
	&	‹dumma›	&	[ˈdŭmˑɑ] 	&	‘much, very, great		\\ \hhline{~}
	&	\hp{ }		&	\hp{ }	&	\hp{`}(amount of, degree of)’		\\
	&	‹yedda›	&	[ˈyĕdˑɑ] {\tl} [ˈyədˑɑ]	&	‘1) still; 2) not yet’		\\
	&	‹afenni›	&	[ɑˈfĕnˑi] {\tl} [ɑˈfənˑi]	&	‘arrogance’		\\
\lspbottomrule
\end{xltabular}\addtocounter{table}{-1}}

Here we simply note the close phonetic similarities in the description
of \tsc{\ili{Alor}-Pantar} languages with descriptions of schwa in
nearby Austronesian languages on Flores as being unstressed,
of short duration, like an epenthetic vowel to break up consonant
clusters, and sometimes triggering lengthening of the following consonant.
We are not claiming contact-induced change in either direction,
but that is something for further investigation.
Since this schwa of short duration often triggering lengthening
of the following consonant was already found in Austronesian
languages to the west
and north, at best its presence in Papuan languages lends support
to maintain it in some Austronesian languages within a limited region.
The phenomena involving schwa in \tsc{\ili{Bima}-Lembata} languages are discussed
in \srf{03-02-01} and \srf{03-03-01}.

%In \srf{02-04-06-01} we discuss various ways people indicate
%schwa /ə/ or other vowels in writing languages with six-vowel
%systems in the region.
%In \srf{02-04-06-02} we discuss particular issues encountered
%in the documentation of Sumba languages,
%along with other transcription issues complicating understanding
%languages in the \tsc{\ili{Bima}-Lembata} subgroup.
%In \srf{02-04-06-03} we discuss particular issues for understanding
%\citeauthor{Onvlee1984}'s (\citeyear{Onvlee1984}) transcriptions,
%necessary to make comparative use of his valuable data source.